% Editable transcription; verify reused text against the original PDF.
\section{Is modern AI (more) psychologically plausible?}

Schank (1980) stated that early AI work [1950's--1970's]
``did not have a very particular psychological flavor'', where
questions about how ``people do various tasks... were not of
particular interest'', although AI researchers were becoming
increasingly interested in cognition in recent years [late 70's].
With the impressive development of AI techniques in the past decades,
we believe that the interest in cognition has only grown.
Although much can be said about the limitations of current AI methods
and models, such as the amount of data required to train state-of-the-art
ANNs and inhuman complexity of some models (such as GPT-3), we see many
ways in which AI has become more psychologically plausible since the
time of Schank's, 1980 article.
We explain this according to David Marr's (1982) three levels of analysis.

\begin{quote}
\small
\textbf{Figure 1 (image omitted).}
Examples of generated images in response to the text prompt
``an armchair in the shape of an avocado'' by DALL-E.
Adapted from \url{https://openai.com/blog/dall-e}.
\end{quote}

In terms of the computational level, the problems modern AI solve
are much more humanlike than the simple toy problems often used
in early AI efforts: modern AI is often used, as witnessed by the
applications introduced before, to solve complex tasks in real-world
contexts and using real-world data.
And these real-world modern AI models sometimes achieve
(super)human performance.
Although most AI models that achieve (super)human performance are
constrained to a single, narrow task, the observation that these
artificial systems perform on par with actual human behavior is
amazing nonetheless.
Interestingly, one direction some researchers take to improve the
performance of their models is to make the task that needs to be
solved more ``human'', or in other words, to align the computational
level between artificial systems and humans.
One characteristic example of this is ``curiosity-based learning'',
a specific implementation of reinforcement learning models which
use ``curiosity'' instead of external reward to motivate learning
(Burda, Edwards, Storkey, \& Klimov, 2018).

Equivalence at the computational level, however, does not necessarily
mean equivalence at the algorithmic level (i.e., a cognitive wheel).
In other words, although modern AI and human cognition aim to solve
similar problems, the way AI and humans implement these solutions
may differ fundamentally.
There is, however, reason to be optimistic in this regard, given
the abundance of studies showing that ANNs and reinforcement
learning models learn representations and contain mechanisms
that have plausible neural correlates.
Especially in the domain of perception, researchers have shown
that the representations learned by ANNs are similar to stimulus
representations in the sensory cortex (e.g., Yamins et al., 2014).
And, importantly, that deeper layers in ANNs contain information
corresponding to stimulus features represented in brain regions
further down the sensory processing hierarchy
(Güçlü \& van Gerven, 2015).
Indeed, at the moment, features from deep neural networks are better
able to explain neural activity during naturalistic vision than
features from classical computational models of the visual cortex
(Khaligh-Razavi \& Kriegeskorte, 2014; Schrimpf et al., 2018).
In addition to ANNs, (deep) reinforcement learning models appear
to have plausible neural correlations, such as the correspondence
of parameters from the temporal difference learning model with
midbrain dopamine activity during animal and human learning
(Schultz, Dayan, \& Montague, 1997; for a review on the correspondence
between reinforcement learning models and the brain, see
O'Doherty, Lee, \& McNamee, 2015).
This correspondence between artificial and biological representations
is further supported by the fact that many AI models, once they are
trained, can often be successfully applied to other tasks and
(perceptual) domains with minimal or no retraining---a technique
called transfer learning (Torrey \& Shavlik, 2010).
The success of transfer learning
(Yosinski, Clune, Bengio, \& Lipson, 2014), which includes the
aforementioned reuse of ANNs in deep reinforcement learning
(Botvinick, Wang, Dabney, Miller, \& Kurth-Nelson, 2020) and deep
language models (Lu, Grover, Abbeel, \& Mordatch, 2021), thus shows
that AI systems can learn generalizable representations,
just like the brain.

Another development that points to AI becoming more human at the
algorithmic level is that researchers started to incorporate concepts
and ideas from psychology and cognitive science in their AI systems.
For example, researchers have shown that implementing an artificial
version of ``attention'' in both language models
(Vaswani et al., 2017) and vision models
(Linsley, Shiebler, Eberhardt, \& Serre, 2018; Xu et al., 2015)
improves their performance (see for a comprehensive review
Lindsay, 2020), which mirrors recent developments in intelligence
research that highlight the role of attention and attentional control
in human intelligence (e.g., Burgoyne \& Engle, 2020).

A similar example in the domain of reinforcement learning is the
implementation of ``episodic memory'' and meta-learning
(``learning to learn'') in artificial reinforcement learning models,
which has been shown to dramatically improve the speed and efficiency
when training these models (Botvinick et al., 2019; Wang, 2021).
Psychologically-infused reinforcement models also appear to have
plausible neural correlations (Wang et al., 2018).

Finally, even though early ANN models were in fact abstractions from
biological neurons and networks, modern AI models are much more human
and brain-like at the implementational level.
Like the brain, modern AI models often feature hierarchical processing
streams that contain a large set of parameters that match the
complexity of the human brain.
These developments include hierarchical layers, convolutions,
synaptic pruning (Blalock, Ortiz, Frankle, \& Guttag, 2020), and
recurrence---all of which have plausible equivalents in the human
brain (Kriegeskorte, 2015).
